\section*{Chapitre \ref{ch:rs:simulation-versus-live} --- Simulation contre production}

\begin{solution}{exo:rs:simulation-versus-live:1}
L'achat en production à 50.00 (10.2~s) correspond à l'achat simulé à 50.00 (10.9~s)~; la vente en production n'a aucune contrepartie simulée, et l'achat simulé à 49.99 n'a aucune contrepartie en production.
100 des 200 actions en production ont été appariées et 100 des 200 actions simulées~: les deux parts valent 50~\%.
\end{solution}

\begin{solution}{exo:rs:simulation-versus-live:2}
Exécution~: $6\,000 \times (20.04 - 20.00) = \$240$. Opportunité~: $4\,000 \times (20.30 - 20.00) = \$1\,200$. Frais~: $6\,000 \times \$0.002 = \$12$. \omterm{def:rs:simulation-versus-live:calibration}{Implementation shortfall}~:
1\,452~\$, dont l'essentiel tient aux actions non achetées.
\end{solution}

\begin{solution}{exo:rs:simulation-versus-live:3}
Les hypothèses interagissent~: l'effet de la latence dépend du modèle de file d'attente et du marché, si bien qu'une estimation d'une étape, les autres étant fixés à leurs valeurs de recherche,
produit une attribution erronée. Un passage ne modifie qu'une seule entrée et conserve la cohérence interne du simulateur, de sorte que les étapes s'additionnent pour donner l'écart total.
\end{solution}

\begin{solution}{exo:rs:simulation-versus-live:4}
L'étape de marché compare deux réalisations du flux d'ordres de l'heure, avec et sans la stratégie~; elles diffèrent par l'impact de la stratégie et par le hasard. Estimer la
part due au hasard à partir de réalisations sans la stratégie (une autre graine du flux d'ordres avec le même prix efficient~: 62.8~\$ par heure ici), et l'impact à partir de nombreuses séances
ou d'une expérience contrôlée (négocier des jours tirés au hasard, chapitre~21).
\end{solution}

\begin{solution}{exo:rs:simulation-versus-live:5}
Les ordres agressifs qui ont rencontré l'ordre de la stratégie et s'y sont arrêtés n'apparaissent plus une fois les messages de la stratégie retirés~; un ordre fantôme ne peut être exécuté que
par les exécutions que contient le rejeu. Dans les séances du chapitre, 68~\% des lots en production provenaient d'ordres agressifs qui ne se sont exécutés que contre le teneur de cotations.
\end{solution}

\begin{solution}{exo:rs:simulation-versus-live:6}
Les autres ordres au prix ont disparu parce que le prix s'éloignait~: annulations de cotations périmées à mesure que le prix efficient se déplaçait, comme dans le mécanisme du simulateur. L'ordre encore
au carnet est exécuté par des intervenants qui savent que le prix a bougé~: sur les séances synthétiques, un mark-out de $-0.85$ tick contre $-0.11$.
\end{solution}

\begin{solution}{exo:rs:simulation-versus-live:7}
\texttt{rs\_simlive.effective\_latency}~: le rejeu exécute 1\,898 lots à latence nulle, moins à chaque latence plus grande, contre 3\,582 en production~: aucune latence ne comble l'écart.
La latence agit sur le moment où les ordres du teneur de cotations arrivent et partent, elle concerne donc les exécutions où le teneur était laissé seul à un prix~; mais, dans le rejeu, ce niveau de prix se vide simplement,
et les transactions qui s'arrêtaient au teneur sont absentes de son entrée, quelle que soit la latence.
\end{solution}

\begin{solution}{exo:rs:simulation-versus-live:8}
Faire coïncider le nombre avec de mauvaises exécutions~: le modèle optimiste ajoute des exécutions favorables qu'un ordre réel n'obtient pas, alors que les exécutions supplémentaires de la production sont les pires. Dans les séances du chapitre,
les exécutions en tête de file (8\,886 lots) dépassent le nombre en production (3\,582) et transforment la perte de 2\,587~\$ observée en production en un profit simulé de 2\,207~\$.
Il faut calibrer sur le P\&L et les mark-outs par groupe d'exécutions, et corriger la cause structurelle (une simulation réactive).
\end{solution}

\begin{solution}{pb:rs:simulation-versus-live:1}
\begin{enumerate}
\item 597 lots par heure en production~; 305 dans le rejeu.
\item 27~\% des lots en production ont été appariés~; 54~\% des lots simulés se sont produits en production.
\item Ce n'est pas une copie bruitée de la production mais un modèle différent~: il produit des exécutions différentes, et non les mêmes avec des erreurs.
\item Données de marché 0.02~s plus une exponentielle de moyenne 0.02~s, transmission des ordres 0.03~s plus une exponentielle de moyenne 0.02~s~; le rejeu de recherche n'en suppose aucune.
\item $-\$50.5$, $-\$56.4$, $-\$112.8$, $-\$431.2$ et $-\$490.9$ par heure.
\item Latence $-\$5.9$, marché $-\$56.4$, exécutions $-\$318.3$, frais $-\$59.7$.
\item Tout cela~: une seconde réalisation de l'heure diffère de la première de 62.8~\$ par heure en moyenne.
\item La latence et les frais sont connus à l'avance~; le marché et les exécutions ont dû être mesurés en production.
\item 68~\%~: ces transactions n'existent qu'avec l'ordre du teneur dans le carnet~; l'entrée du rejeu, c'est-à-dire le marché sans cet ordre, ne les contient pas.
\item 21~\%, avec un mark-out à 10~secondes de $-0.85$ tick contre $-0.11$ pour les autres exécutions en production.
\item Son annulation est en transit lorsque les autres ordres quittent le prix, ce qu'ils font parce que le prix s'éloigne.
\item Annuler lorsque les autres ordres à son prix s'amenuisent, ou coter un tick derrière le meilleur cours lorsque la file est courte.
\item \textbf{Résultat nommé.} Du $-\$50.5$ par heure du rejeu de recherche au $-\$490.9$ de la production~: latence $-\$5.9$, marché $-\$56.4$ (dans la marge du hasard), exécutions en production
$-\$318.3$, frais $-\$59.7$. L'écart se loge dans les exécutions~: 68~\% des lots en production proviennent de transactions qui s'arrêtaient au teneur, et 21~\% ont été prises alors qu'il était seul à son prix.
\item Rien~: aucune latence ne fait exécuter au rejeu autant de lots qu'en production.
\item Il exécute plus de lots que la production (8\,886 contre 3\,582 sur les séances) et affiche un profit (2\,207~\$ contre une perte en production de 2\,587~\$).
\item Une simulation réactive dans laquelle les ordres de la stratégie font partie du marché, pour des stratégies qui représentent une part notable du flux à leurs prix.
\item La parité dans les deux sens, les étapes de la décomposition, les mark-outs par groupe d'exécutions, et l'\omterm{def:rs:simulation-versus-live:calibration}{implementation shortfall} de chaque ordre parent.
\item Prendre le prix de décision de chaque débouclage (au moment où la limite d'inventaire est atteinte) comme \omterm{def:rs:simulation-versus-live:calibration}{prix d'arrivée}, et décomposer son coût en exécution, coût d'opportunité (la part non débouclée) et frais.
\item Moins de transactions s'arrêtent au teneur et moins de niveaux lui sont laissés seul~: les angles morts du rejeu se réduisent, et le rejeu devient une meilleure estimation.
\item Être rapproché de la production chaque jour, les chiffres du rapprochement décidant de ce qui change dans le modèle.
\end{enumerate}
\end{solution}

\begin{solution}{iq:rs:simulation-versus-live:1}
Apparier d'abord les exécutions (parité dans les deux sens)~: des exécutions manquantes ou en trop désignent le \omterm{def:rs:event-driven-backtests:fill}{modèle d'exécution} et la latence~; décomposer ensuite l'écart de P\&L par passages (latence, marché, exécutions, frais, coûts),
puis examiner les mark-outs des exécutions non appariées. Vérifier aussi les éléments simples~: frais, emprunt, opérations sur titres, synchronisation des horloges.
\end{solution}

\begin{solution}{iq:rs:simulation-versus-live:2}
L'écart entre le P\&L d'un portefeuille papier négocié en totalité aux prix de décision et le P\&L réalisé (Perold, 1988). Le décomposer en retard (de la décision à l'arrivée),
coût d'exécution de la part exécutée par rapport au \omterm{def:rs:simulation-versus-live:calibration}{prix d'arrivée}, coût d'opportunité de la part non exécutée, et frais.
\end{solution}

\begin{solution}{iq:rs:simulation-versus-live:3}
Rejouer une séance de production dans le simulateur avec les ordres réels et leurs instants effectifs, puis comparer les exécutions~: la part des exécutions de production reproduites et la part des exécutions simulées
qui se sont produites. Cela permet de détecter des erreurs dans la logique d'exécution, des décalages d'horloge, de mauvaises hypothèses de latence et des angles morts structurels (comme les transactions qui n'existent que du fait de la
stratégie).
\end{solution}

\begin{solution}{iq:rs:simulation-versus-live:4}
Le marché enregistré contient les conséquences des ordres de la stratégie~: des transactions qui se sont arrêtées sur ces ordres, les réactions des autres intervenants. Retirer les messages de la stratégie laisse un marché qui n'a jamais existé
et n'aurait pas existé sans elle~; les conserver compte deux fois la liquidité. Dans les séances du chapitre, 68~\% des lots en production provenaient de transactions absentes du
marché sans le teneur de cotations.
\end{solution}

\begin{solution}{iq:rs:simulation-versus-live:5}
Ajuster peu de paramètres, à signification physique (latence, paramètre d'un modèle de file d'attente), sur une période et les vérifier sur une autre~; juger sur le P\&L et les mark-outs par groupe d'exécutions, ainsi que sur les nombres~;
refuser les paramètres qui font coïncider les nombres en ajoutant de mauvaises exécutions~; privilégier enfin un changement structurel lorsque la parité montre des exécutions différentes plutôt que bruitées.
\end{solution}

\begin{solution}{iq:rs:simulation-versus-live:6}
Faire tourner en \omterm{def:rs:simulation-versus-live:reconciliation}{trading fantôme} le simulateur de chaque stratégie sur le flux de production~; chaque nuit, pour chaque stratégie, calculer la parité, les étapes de la décomposition, les mark-outs par groupe d'exécutions et l'\omterm{def:rs:simulation-versus-live:calibration}{implementation
shortfall}~; signaler les dérives au-delà du bruit propre à chaque stratégie~; conserver l'ensemble avec la version du simulateur, afin qu'un changement du simulateur puisse être contrôlé par rapport à son
historique (chapitre~29).
\end{solution}
